\documentclass{article}
\begin{document}
\section{Sources}
The method goes back to \cite{gauss} and was refined in \cite{jacobi,seidel}; see also
\cite[p.~12]{gauss} for the original statement, and \cite[Chapter~3]{seidel,young}.
A typesetting system takes a description of a document and decides where every character goes. It breaks paragraphs into lines, stacks the lines into pages and places headings, notes and figures where the reader expects them. Each of these steps follows rules that can be written down precisely, and a good program follows them exactly.
\begin{thebibliography}{9}
\bibitem{gauss} C. F. Gauss. Letter to Gerling, 1823.
\bibitem{jacobi} C. G. J. Jacobi. On a new way of solving linear equations. Astronomische Nachrichten 22, 1845.
\bibitem{seidel} L. Seidel. On a process for solving the equations of least squares by successive approximation. Abhandlungen der Bayerischen Akademie 11, 1874.
\bibitem{young} D. M. Young. Iterative methods for solving partial difference equations of elliptic type. Thesis, 1950.
\end{thebibliography}
\end{document}
